% ---------------------------------------------------------------------
% STATUT : SQUELETTE - c'est le chapitre NEUF et probablement la colonne
%          vertébrale du mémoire. Contenu arrêté avec Hugo le 21/07/2026.
% EN ATTENTE : confirmation du paradigme d'ensemble (verdict du panel).
% À FAIRE AVANT RÉDACTION : coder le cran 1 (u_ij = a_i b_j) et produire
%          la figure de l'échelle de repli (SCR et classement par cran).
% ---------------------------------------------------------------------
\chapter{Décomposition de la propagation et échelle de repli}
\label{ch:decomposition}

\begin{cle}
\textbf{Thèse du chapitre.} La probabilité de propagation se décompose en une part que la
donnée fixe et une part qu'elle ne fixe pas. On pose
\[
  p_{ij} \;=\; p_j \;+\; u_{ij},
\]
où $p_j$ est l'\textbf{attractivité de la cible} $j$ (à quel point $j$ est touché, quelle que
soit la provenance) et $u_{ij}$ l'\textbf{excès propre à la source} $i$. La première est une
marginale, donc calibrable ; la seconde \emph{est} la direction, donc la contribution du
mémoire, et c'est elle qui n'est pas identifiable. Le chapitre attaque l'estimation la plus
ambitieuse de $u_{ij}$, montre où elle casse, et \textbf{recule par crans} jusqu'au modèle
nul $u \equiv 0$.
\end{cle}

\section{La décomposition et sa normalisation}

% À RÉDIGER. Points obligatoires :
%  - définir p_ij dans le modèle de cascade (lien avec W_jk du chapitre 4)
%  - p_j = attractivité de la cible ; u_ij = excès propre à la source
%  - PIÈGE 1 : la décomposition n'est PAS identifiée sans normalisation.
%    Pour toute constante c_j, (p_j + c_j, u_ij - c_j) donne le même p_ij.
%    Remède : imposer  moyenne_i u_ij = 0  pour chaque j.
%    (problème classique des effets à deux voies)
%  - PIÈGE 2 : la forme additive peut sortir de [0,1].
%    Remède : contraindre u_ij dans [-p_j, 1-p_j], ou passer en logit :
%    logit p_ij = logit p_j + u_ij  (recommandé)
%  - traiter le cas de la diagonale u_jj
%  - relier au rayon spectral : la contrainte rho(W) < 1 borne u

\section{L'asymétrie du mécanisme n'est pas l'asymétrie des taux}

% À RÉDIGER. Point clé, souvent mal compris :
%  sous u = 0 on a p_ij = p_j et p_ji = p_i, qui diffèrent en général.
%  Cette asymétrie-là ne vient QUE de l'attractivité des cibles, ce n'est
%  pas un mécanisme directionnel. Le vrai mécanisme est u_ij - u_ji.
%  Conséquence : le test de knockout d'asymétrie teste en réalité u.

\section{Le test : la source compte-t-elle au-delà de la cible ?}

% À RÉDIGER. Reformuler le résultat existant comme un test d'hypothèse :
%     H0 : u_ij = 0 pour tout i,j
%  Le placebo par permutation donne z = -0,33 : H0 n'est PAS rejetée sur
%  les données publiques. Ce n'est plus « je n'ai pas su calibrer W »,
%  c'est « la donnée ne permet pas de distinguer la source de la cible ».
%  Renvoyer au chapitre \ref{ch:identifiabilite} pour les deux sources.

\section{L'échelle de repli}

% À RÉDIGER + À CODER. Les cinq crans, du plus dur au repli :
%   Cran 0 : u_ij libre (20 paramètres hors diagonale). Non identifiable.
%   Cran 1 : u_ij = a_i b_j (rang 1 : émetteurs a_i, récepteurs b_j),
%            environ 10 paramètres.  <-- JAMAIS TESTÉ, à coder en premier
%   Cran 2 : u_ij non estimé mais BORNÉ (identification partielle),
%            SCR reporté en bornes sur l'ensemble admissible
%   Cran 3 : u_ij élicité (protocole de Cooke, chapitre \ref{ch:elicitation})
%   Cran 4 : u_ij = 0, soit p_ij = p_j. MODÈLE NUL, entièrement identifié,
%            mais la dépendance à l'ordre y disparaît.
%
% LIVRABLE CENTRAL : une figure unique donnant, pour chaque cran, le SCR
% et le classement des piliers. Toute la contribution du mémoire se lit
% comme l'écart au cran 4.

\begin{cle}
\textbf{Pourquoi cette construction tient devant un jury.} Elle applique la méthode demandée
par le tuteur (attaquer le modèle le plus ambitieux, reculer par crans documentés) et elle
coïncide avec un paradigme statistique reconnu, l'\textbf{identification partielle} : on
n'estime pas un paramètre que la donnée ne détermine pas, on caractérise l'ensemble de ses
valeurs admissibles et l'on montre ce que le capital devient sur cet ensemble. La limite
cesse d'être une excuse pour devenir la méthode.
\end{cle}
